\chapter{Chapter06 — Small Bias Fourier Algorithm}
%%% generated by the blueprint pipeline from FABL.Chapter06.LearningAndTesting.SmallBiasFourierAlgorithm
\section{Overview}

Declarations of the Lean module \texttt{FABL.Chapter06.LearningAndTesting.SmallBiasFourierAlgorithm}.

\section{Declarations}

\begin{theorem}[Uniform pushforward under reindexing]
\lean{FABL.ProbabilityDensity.IsBiased.uniformPushforward_comp_equiv}
\label{FABL.ProbabilityDensity.IsBiased.uniformPushforward_comp_equiv}
\leanok
\uses{FABL.F₂Cube, FABL.ProbabilityDensity.IsBiased, FABL.ProbabilityDensity.expectation_uniformPushforward, FABL.ProbabilityDensity.isBiased_iff_expectation, FABL.ProbabilityDensity.uniformPushforward, FABL.vectorWalshCharacter}
\difficulty{4}
Let $\Omega$ and $\Omega'$ be finite nonempty sets, let $g$ map $\Omega$ to the Boolean
cube $\mathbb{F}_2^n$, and let $\varepsilon\in\mathbb{R}$. If the distribution on the
cube induced by a uniform element of $\Omega$ is $\varepsilon$-biased, then the
distribution induced by $g\circ e$ from a uniform element of $\Omega'$ is also
$\varepsilon$-biased, for every bijection $e:\Omega'\to\Omega$.
\end{theorem}

\begin{definition}[Small-bias Fourier input]
\lean{FABL.SmallBiasFourierInput}
\statementsource{Analysis-of-Boolean-Functions-by-Ryan-ODonnell}{
  pdf-ffbe9c7b3306-p153-b002,
  pdf-ffbe9c7b3306-p153-b003,
  pdf-ffbe9c7b3306-p157-b008,
  pdf-ffbe9c7b3306-p157-b009
}
\label{FABL.SmallBiasFourierInput}
\leanok
\uses{FABL.SmallBiasInput}
A small-bias Fourier input consists of a small-bias input together with a natural-number
bound $s$ satisfying $1\le s$, where $s$ is an integral upper bound for the target's
Fourier $1$-norm.
\end{definition}

\begin{definition}[Fourier input bias]
\lean{FABL.SmallBiasFourierInput.epsilon}
\statementsource{Analysis-of-Boolean-Functions-by-Ryan-ODonnell}{
  pdf-ffbe9c7b3306-p153-b002,
  pdf-ffbe9c7b3306-p157-b008
}
\label{FABL.SmallBiasFourierInput.epsilon}
\leanok
\uses{FABL.SmallBiasFourierInput, FABL.SmallBiasInput.epsilon}
For a small-bias Fourier input whose bias has numerator $a$ and denominator $b$, its
real bias parameter is $a/b$.
\end{definition}

\begin{definition}[Generator input parameters]
\lean{FABL.SmallBiasFourierInput.generatorInput}
\statementsource{Analysis-of-Boolean-Functions-by-Ryan-ODonnell}{
  pdf-ffbe9c7b3306-p153-b002,
  pdf-ffbe9c7b3306-p153-b003
}
\label{FABL.SmallBiasFourierInput.generatorInput}
\leanok
\uses{FABL.SmallBiasFourierInput, FABL.SmallBiasInput}
Given a small-bias Fourier input with bias numerator $a$, denominator $b$, and Fourier
bound $B$, its associated generator input has output dimension $n$, numerator $a$, and
denominator $bB$. It retains the hypotheses $n>0$, $a>0$, and $2a\le bB$.
\end{definition}

\begin{theorem}[Generator bias scaling]
\lean{FABL.SmallBiasFourierInput.generatorInput_epsilon}
\label{FABL.SmallBiasFourierInput.generatorInput_epsilon}
\leanok
\uses{FABL.SmallBiasFourierInput, FABL.SmallBiasFourierInput.epsilon, FABL.SmallBiasFourierInput.generatorInput, FABL.SmallBiasInput.epsilon}
\difficulty{4}
For a small-bias Fourier input with target bias $\varepsilon$ and Fourier bound $B$, the bias
parameter of its associated generator input equals the target bias divided by $B$.
\end{theorem}

\begin{definition}[Sample count]
\lean{FABL.SmallBiasFourierInput.sampleCount}
\statementsource{Analysis-of-Boolean-Functions-by-Ryan-ODonnell}{
  pdf-ffbe9c7b3306-p153-b002,
  pdf-ffbe9c7b3306-p153-b003
}
\label{FABL.SmallBiasFourierInput.sampleCount}
\leanok
\uses{FABL.SmallBiasFourierInput, FABL.SmallBiasFourierInput.generatorInput, FABL.SmallBiasInput.fieldDegree}
For a small-bias Fourier input, let $k$ be the field degree of its associated generator
input. Its sample count is $2^k\cdot 2^k$.
\end{definition}

\begin{theorem}[Positivity of the sample count]
\lean{FABL.SmallBiasFourierInput.sampleCount_pos}
\label{FABL.SmallBiasFourierInput.sampleCount_pos}
\leanok
\uses{FABL.SmallBiasFourierInput, FABL.SmallBiasFourierInput.sampleCount}
\difficulty{1}
For every small-bias Fourier input, its sample count is positive.
\end{theorem}

\begin{definition}[Deterministic construction]
\lean{FABL.SmallBiasFourierInput.construction}
\statementsource{Analysis-of-Boolean-Functions-by-Ryan-ODonnell}{
  pdf-ffbe9c7b3306-p153-b002,
  pdf-ffbe9c7b3306-p153-b003
}
\label{FABL.SmallBiasFourierInput.construction}
\leanok
\uses{FABL.ExecutableSmallBiasConstruction, FABL.SmallBiasFourierInput, FABL.SmallBiasFourierInput.generatorInput, FABL.deterministicSmallBiasAlgorithm}
For a small-bias Fourier input, its construction is the deterministic small-bias
algorithm applied to the associated generator input. The resulting executable
construction has output dimension equal to that input's output dimension.
\end{definition}

\begin{definition}[Generated sample map]
\lean{FABL.SmallBiasFourierInput.sample}
\statementsource{Analysis-of-Boolean-Functions-by-Ryan-ODonnell}{
  pdf-ffbe9c7b3306-p153-b003,
  pdf-ffbe9c7b3306-p157-b009
}
\label{FABL.SmallBiasFourierInput.sample}
\leanok
\uses{FABL.F₂Cube, FABL.SmallBiasFourierInput, FABL.SmallBiasFourierInput.generatorInput, FABL.SmallBiasFourierInput.sampleCount, FABL.SmallBiasInput.fieldDegree, FABL.SmallBiasInput.fieldDegree_pos, FABL.binaryVectorPairEnumerationEquiv, FABL.deterministicSmallBiasAlgorithm, FABL.executableSmallBiasGenerator}
For a small-bias Fourier input, the sample is a map from the indices $0,\ldots,N-1$ to
binary vectors of length $n$, where $N$ is its sample count and $n$ is the associated
generator input's output dimension. At each index, the map enumerates the corresponding
ordered pair of binary seeds of length $k$, where $k$ is the generator input's field
degree, and applies the executable small-bias generator using the irreducible-polynomial
field implementation selected by the deterministic algorithm.
\end{definition}

\begin{definition}[Sample probability density]
\lean{FABL.SmallBiasFourierInput.sampleDensity}
\statementsource{Analysis-of-Boolean-Functions-by-Ryan-ODonnell}{
  pdf-ffbe9c7b3306-p153-b002,
  pdf-ffbe9c7b3306-p157-b009
}
\label{FABL.SmallBiasFourierInput.sampleDensity}
\leanok
\uses{FABL.ProbabilityDensity, FABL.ProbabilityDensity.uniformPushforward, FABL.SmallBiasFourierInput, FABL.SmallBiasFourierInput.generatorInput, FABL.SmallBiasFourierInput.sample, FABL.SmallBiasFourierInput.sampleCount, FABL.SmallBiasFourierInput.sampleCount_pos}
For a small-bias Fourier input, its sample density is the uniform pushforward of the
uniform distribution on the finite sample index set under its sample map. It is a
probability density on the binary cube of dimension equal to the associated generator
input's output dimension.
\end{definition}

\begin{theorem}[Sample is biased]
\lean{FABL.SmallBiasFourierInput.sample_isBiased}
\statementsource{Analysis-of-Boolean-Functions-by-Ryan-ODonnell}{
  pdf-ffbe9c7b3306-p153-b003,
  pdf-ffbe9c7b3306-p153-b004
}
\label{FABL.SmallBiasFourierInput.sample_isBiased}
\leanok
\uses{FABL.F₂Cube, FABL.ProbabilityDensity.IsBiased, FABL.ProbabilityDensity.IsBiased.uniformPushforward_comp_equiv, FABL.ProbabilityDensity.uniformPushforward, FABL.SmallBiasFourierInput, FABL.SmallBiasFourierInput.epsilon, FABL.SmallBiasFourierInput.generatorInput, FABL.SmallBiasFourierInput.generatorInput_epsilon, FABL.SmallBiasFourierInput.sampleCount, FABL.SmallBiasFourierInput.sampleCount_pos, FABL.SmallBiasFourierInput.sampleDensity, FABL.SmallBiasInput.epsilon, FABL.SmallBiasInput.fieldDegree, FABL.SmallBiasInput.fieldDegree_pos, FABL.binaryVectorPairEnumerationEquiv, FABL.deterministicSmallBiasAlgorithm, FABL.deterministicSmallBiasAlgorithm_spec, FABL.executableSmallBiasGenerator, FABL.executableSmallBiasGeneratorDensity}
\difficulty{4}
% dataset-pipeline: stub -- prose pending build_informal_statements.py
Statement pending regeneration.
\end{theorem}

\begin{definition}[Deterministic estimator cost]
\lean{FABL.deterministicSmallBiasFourierEstimatorCost}
\statementsource{Analysis-of-Boolean-Functions-by-Ryan-ODonnell}{pdf-ffbe9c7b3306-p157-b008}
\label{FABL.deterministicSmallBiasFourierEstimatorCost}
\leanok
\uses{FABL.LearningCost, FABL.SmallBiasFourierInput, FABL.SmallBiasFourierInput.generatorInput, FABL.SmallBiasFourierInput.sampleCount, FABL.deterministicSmallBiasWork, FABL.smallBiasFourierEstimatorCost}
Given a small-bias Fourier input and a finite set $U$ of output coordinates, the
deterministic Fourier estimator cost is the sum of the cost of the deterministic
generator's local work and the Fourier estimator's cost on $U$ using the prescribed
sample count. Its random-example and query counts are zero, apart from any contributions
in the estimator cost.
\end{definition}

\begin{definition}[Deterministic Fourier estimator program]
\lean{FABL.deterministicSmallBiasFourierEstimatorProgram}
\statementsource{Analysis-of-Boolean-Functions-by-Ryan-ODonnell}{
  pdf-ffbe9c7b3306-p157-b008,
  pdf-ffbe9c7b3306-p157-b009
}
\label{FABL.deterministicSmallBiasFourierEstimatorProgram}
\leanok
\uses{FABL.DeterministicQueryProgram, FABL.SmallBiasFourierInput, FABL.SmallBiasFourierInput.generatorInput, FABL.SmallBiasFourierInput.sample, FABL.deterministicSmallBiasWork, FABL.smallBiasFourierEstimatorProgram}
Given a small-bias Fourier input and a finite set $U$ of output coordinates, the
deterministic Fourier estimator program first charges the local computation work of the
deterministic generator, then runs the Fourier estimator program on the generated sample
and $U$. Its query type is the set of sign vectors on the output cube, and its answer
and output types are both the real numbers.
\end{definition}

\begin{theorem}[Cost of the deterministic small-bias Fourier estimator]
\lean{FABL.DeterministicQueryProgram.runWithCost_deterministicSmallBiasFourierEstimatorProgram}
\label{FABL.DeterministicQueryProgram.runWithCost_deterministicSmallBiasFourierEstimatorProgram}
\leanok
\uses{FABL.DeterministicQueryProgram.runWithCost, FABL.SmallBiasFourierInput, FABL.SmallBiasFourierInput.generatorInput, FABL.SmallBiasFourierInput.sample, FABL.deterministicSmallBiasFourierEstimatorCost, FABL.deterministicSmallBiasFourierEstimatorProgram, FABL.smallBiasFourierEstimate}
\difficulty{1}
Let the input consist of a positive dimension $n$, positive integers $a$ and $b$ with
$2a \le b$, and an integer Fourier bound $B \ge 1$. Let $f$ be a real-valued function on
the sign cube of dimension $n$, and let $U$ be a finite set of coordinates in that cube.
Running the deterministic query program with oracle $f$ returns the average of
$f(x)\prod_{i\in U}x_i$ over the program's prescribed sample, together with its cost:
zero random examples and membership queries, and local work equal to the deterministic
construction work plus the Fourier estimator work.
\end{theorem}

\begin{theorem}[Abs deterministic small bias fourier estimate sub fourier coeff le]
\lean{FABL.abs_deterministicSmallBiasFourierEstimate_sub_fourierCoeff_le}
\statementsource{Analysis-of-Boolean-Functions-by-Ryan-ODonnell}{
  pdf-ffbe9c7b3306-p157-b003,
  pdf-ffbe9c7b3306-p157-b004,
  pdf-ffbe9c7b3306-p157-b009
}
\label{FABL.abs_deterministicSmallBiasFourierEstimate_sub_fourierCoeff_le}
\leanok
\uses{FABL.SmallBiasFourierInput, FABL.SmallBiasFourierInput.epsilon, FABL.SmallBiasFourierInput.generatorInput, FABL.SmallBiasFourierInput.sample, FABL.SmallBiasFourierInput.sampleCount, FABL.SmallBiasFourierInput.sampleCount_pos, FABL.SmallBiasFourierInput.sampleDensity, FABL.SmallBiasFourierInput.sample_isBiased, FABL.SmallBiasInput.epsilon_le_half, FABL.SmallBiasInput.epsilon_pos, FABL.abs_smallBiasFourierEstimate_sub_fourierCoeff_le_parameter, FABL.fourierCoeff, FABL.fourierOneNorm, FABL.smallBiasFourierEstimate}
\difficulty{3}
% dataset-pipeline: stub -- prose pending build_informal_statements.py
Statement pending regeneration.
\end{theorem}

\begin{theorem}[Deterministic small bias fourier estimator cost queries]
\lean{FABL.deterministicSmallBiasFourierEstimatorCost_queries}
\statementsource{Analysis-of-Boolean-Functions-by-Ryan-ODonnell}{
  pdf-ffbe9c7b3306-p157-b008,
  pdf-ffbe9c7b3306-p157-b009
}
\label{FABL.deterministicSmallBiasFourierEstimatorCost_queries}
\leanok
\uses{FABL.SmallBiasFourierInput, FABL.SmallBiasFourierInput.generatorInput, FABL.SmallBiasFourierInput.sampleCount, FABL.deterministicSmallBiasFourierEstimatorCost}
\difficulty{2}
% dataset-pipeline: stub -- prose pending build_informal_statements.py
Statement pending regeneration.
\end{theorem}

\begin{theorem}[No random-example cost]
\lean{FABL.deterministicSmallBiasFourierEstimatorCost_randomExamples}
\label{FABL.deterministicSmallBiasFourierEstimatorCost_randomExamples}
\leanok
\uses{FABL.SmallBiasFourierInput, FABL.SmallBiasFourierInput.generatorInput, FABL.deterministicSmallBiasFourierEstimatorCost}
\difficulty{1}
For a small-bias Fourier estimation input whose accuracy data satisfy $n>0$, positive
numerator, and twice the numerator at most the denominator, and whose Fourier bound is
at least $1$, the estimator's random-example cost is zero for every finite coordinate
set $U$ in the input dimension.
\end{theorem}

\begin{definition}[Polynomial budget]
\lean{FABL.SmallBiasFourierInput.polynomialBudget}
\statementsource{Analysis-of-Boolean-Functions-by-Ryan-ODonnell}{
  pdf-ffbe9c7b3306-p153-b002,
  pdf-ffbe9c7b3306-p157-b008
}
\label{FABL.SmallBiasFourierInput.polynomialBudget}
\leanok
\uses{FABL.SmallBiasFourierInput, FABL.SmallBiasFourierInput.generatorInput, FABL.SmallBiasInput.scale}
For a small-bias Fourier input, the polynomial budget is defined as $2^{18}(s+1)^8$,
where $s$ is the scale of its generator input.
\end{definition}

\begin{theorem}[Sample count upper bound]
\lean{FABL.SmallBiasFourierInput.sampleCount_le_four_mul_scale_sq}
\label{FABL.SmallBiasFourierInput.sampleCount_le_four_mul_scale_sq}
\leanok
\uses{FABL.SmallBiasFourierInput, FABL.SmallBiasFourierInput.generatorInput, FABL.SmallBiasFourierInput.sampleCount, FABL.SmallBiasInput.fieldDegree, FABL.SmallBiasInput.fieldSize_le_two_scale, FABL.SmallBiasInput.scale}
\difficulty{3}
For a small-bias Fourier input, let $s$ be the scale of its associated generator input.
Its sample count is at most $4s^2$.
\end{theorem}

\begin{theorem}[Work of the Fourier estimator]
\lean{FABL.deterministicSmallBiasFourierEstimatorCost_work}
\label{FABL.deterministicSmallBiasFourierEstimatorCost_work}
\leanok
\uses{FABL.SmallBiasFourierInput, FABL.SmallBiasFourierInput.generatorInput, FABL.SmallBiasFourierInput.sampleCount, FABL.deterministicSmallBiasFourierEstimatorCost, FABL.deterministicSmallBiasWork, FABL.smallBiasFourierEstimatorWork}
\difficulty{2}
For a small-bias Fourier estimation input with $n>0$, positive accuracy numerator, twice
the numerator at most the denominator, and Fourier bound at least $1$, and for any
finite coordinate set $U$ in the input dimension, the estimator's work is the sum of the
deterministic small-bias construction work, the sample count, and the Fourier estimator
work at that sample count and set $U$.
\end{theorem}

\begin{theorem}[Deterministic small bias fourier estimator cost resource bounds]
\lean{FABL.deterministicSmallBiasFourierEstimatorCost_resource_bounds}
\statementsource{Analysis-of-Boolean-Functions-by-Ryan-ODonnell}{pdf-ffbe9c7b3306-p157-b008}
\label{FABL.deterministicSmallBiasFourierEstimatorCost_resource_bounds}
\leanok
\uses{FABL.SmallBiasFourierInput, FABL.SmallBiasFourierInput.generatorInput, FABL.SmallBiasFourierInput.polynomialBudget, FABL.SmallBiasFourierInput.sampleCount, FABL.SmallBiasFourierInput.sampleCount_le_four_mul_scale_sq, FABL.SmallBiasInput.n_le_scale, FABL.SmallBiasInput.polynomialBudget, FABL.SmallBiasInput.scale, FABL.deterministicSmallBiasFourierEstimatorCost, FABL.deterministicSmallBiasFourierEstimatorCost_queries, FABL.deterministicSmallBiasFourierEstimatorCost_randomExamples, FABL.deterministicSmallBiasFourierEstimatorCost_work, FABL.deterministicSmallBiasWork, FABL.deterministicSmallBiasWork_le_polynomialBudget, FABL.smallBiasFourierEstimatorWork}
\difficulty{5}
% dataset-pipeline: stub -- prose pending build_informal_statements.py
Statement pending regeneration.
\end{theorem}

\begin{definition}[Small-bias Fourier task]
\lean{FABL.SmallBiasFourierTask}
\statementsource{Analysis-of-Boolean-Functions-by-Ryan-ODonnell}{pdf-ffbe9c7b3306-p157-b008}
\label{FABL.SmallBiasFourierTask}
\leanok
\uses{FABL.SmallBiasFourierInput, FABL.SmallBiasFourierInput.generatorInput}
A small-bias Fourier task is a pair consisting of a small-bias Fourier input and a
finite set of coordinates of the dimension of its generator input.
\end{definition}

\begin{definition}[Fourier task scale]
\lean{FABL.smallBiasFourierTaskScale}
\label{FABL.smallBiasFourierTaskScale}
\leanok
\uses{FABL.SmallBiasFourierInput.generatorInput, FABL.SmallBiasFourierTask, FABL.SmallBiasInput.scale}
A small-bias Fourier task consists of a small-bias input together with a finite set of
coordinates in its generator dimension. Its scale is the scale parameter of the
generator input determined by the small-bias input and its Fourier bound.
\end{definition}

\begin{theorem}[Deterministic small bias fourier estimator queries is big o]
\lean{FABL.deterministicSmallBiasFourierEstimator_queries_isBigO}
\statementsource{Analysis-of-Boolean-Functions-by-Ryan-ODonnell}{
  pdf-ffbe9c7b3306-p157-b008,
  pdf-ffbe9c7b3306-p157-b009
}
\label{FABL.deterministicSmallBiasFourierEstimator_queries_isBigO}
\leanok
\uses{FABL.SmallBiasFourierTask, FABL.deterministicSmallBiasFourierEstimatorCost, FABL.deterministicSmallBiasFourierEstimatorCost_resource_bounds, FABL.smallBiasFourierTaskScale}
\difficulty{3}
% dataset-pipeline: stub -- prose pending build_informal_statements.py
Statement pending regeneration.
\end{theorem}

\begin{theorem}[Asymptotic work of the Fourier estimator]
\lean{FABL.deterministicSmallBiasFourierEstimator_work_isBigO}
\label{FABL.deterministicSmallBiasFourierEstimator_work_isBigO}
\leanok
\uses{FABL.SmallBiasFourierTask, FABL.deterministicSmallBiasFourierEstimatorCost, FABL.deterministicSmallBiasFourierEstimatorCost_resource_bounds, FABL.smallBiasFourierTaskScale}
\difficulty{3}
As the scale of small-bias Fourier estimation tasks tends to infinity, the estimator's
work is $O((\text{scale}+1)^8)$. Here each task consists of accuracy data with positive
dimension, positive numerator, and twice the numerator at most the denominator, a
Fourier bound of at least $1$, and a finite coordinate set in the resulting dimension.
\end{theorem}
